\chapter{Introdução}
\label{Introdução}

Nos últimos anos, os LLMs têm sido amplamente utilizados…

No contexto educacional e profissional…

Entre esses contextos, certificações como o PMP se destacam…

Entretanto, o uso dessas ferramentas levanta questões sobre confiabilidade…

Diante disso, este trabalho investiga…



\section{Contexto e Motivação}

Ideia central: situar o problema no mundo real.

O que colocar:

Crescente uso de LLMs (ChatGPT, Gemini, etc.) como apoio a estudo e certificações.

Importância da certificação PMP no mercado (gestão de projetos, carreira, padronização).

Mudança recente do exame (ECO 2026 → mais cenário, menos decoreba).

Característica do PMP:
não é só factual → exige julgamento situacional
mistura predictive, agile e hybrid
Motivação prática: candidatos já estão usando LLMs → mas não sabemos se podem confiar.

Fechamento da seção:

“Apesar da crescente adoção, a confiabilidade dessas ferramentas no contexto do PMP atual ainda não está bem compreendida.”



\section{Problema de Pesquisa}

Ideia central: mostrar que o problema ainda não foi resolvido pela literatura.

Estruture em 3 camadas:

(i) O que já existe
Estudos de LLMs em:
exames acadêmicos
certificações técnicas
contexto educacional
Estudos em Agile/Scrum 


(ii) Limitação desses estudos
focam em:
acurácia geral
perguntas factuais
pouco foco em:
domínios híbridos (predictive + agile)
questões situacionais
certificações profissionais complexas

(iii) Gap explícito

“Há uma ausência de estudos empíricos que avaliem sistematicamente o desempenho de LLMs em questões do exame PMP atual, considerando diferentes estratégias de prompting, tipos de questão e domínios de conhecimento.”




\section{Objetivo}

Direto, uma frase bem técnica:

Avaliar o desempenho de Modelos de Linguagem de Grande Porte na resolução de questões alinhadas ao exame PMP atual, investigando o impacto de diferentes estratégias de prompting na acurácia, estabilidade e padrões de erro.

Listar os Objetivos Específicos

Curto, operacional, alinhado à metodologia:

Construir/curar um dataset de questões alinhadas ao PMP (ECO 2026)
Avaliar o desempenho de LLMs sob diferentes estratégias de prompting
Analisar o desempenho por:
domínio (People, Process, Business Environment)
tipo de questão
Identificar padrões de erro e limitações
Derivar implicações para uso de LLMs na preparação para certificação

Evite objetivos genéricos tipo “analisar” sem dizer como.




\section{Metodologia}
Aqui NÃO é detalhamento completo (isso vai no Capítulo 4).
É só um overview do desenho do estudo.

Inclua:

tipo de estudo:

estudo empírico quantitativo com apoio de análise qualitativa

objeto:

questões alinhadas ao PMP atual

abordagem:

avaliação de LLMs sob diferentes estratégias de prompting

variáveis analisadas:
acurácia
desempenho por domínio
desempenho por tipo de questão
padrões de erro


procedimento (alto nível):
execução isolada das questões
coleta estruturada das respostas
análise comparativa




\section{Contribuições Esperadas}

Aqui é onde muitas introduções ficam fracas — não deixe genérico.

Separe em 3 níveis:

Científica
evidência empírica sobre LLMs em certificação profissional

Metodológica
protocolo de avaliação reprodutível (dataset + prompts + métricas)

Prática
diretrizes para uso de LLMs por:
candidatos PMP
cursos preparatórios
instrutores


\section{Organização do Documento}

Este documento está estruturado em seis capítulos, organizados de forma a conduzir o leitor desde a contextualização do problema até a apresentação dos resultados e conclusões da pesquisa.

O Capítulo~\ref{sec:fundamentacao} apresenta a fundamentação teórica, abordando os conceitos essenciais relacionados à certificação PMP, aos Modelos de Linguagem de Grande Porte (LLMs) e às estratégias de \textit{prompt engineering}, que sustentam a proposta deste trabalho.

O Capítulo~\ref{sec:relacionados} discute os trabalhos relacionados, situando esta pesquisa no contexto do estado da arte e destacando as principais lacunas existentes na literatura.

O Capítulo~\ref{sec:metodologia} descreve a metodologia adotada, detalhando o desenho experimental, o dataset utilizado, as estratégias de prompting e as métricas de avaliação consideradas.

O Capítulo~\ref{sec:resultados} apresenta os resultados preliminares obtidos, incluindo análises quantitativas e qualitativas que evidenciam tendências iniciais e a viabilidade da abordagem proposta.

Por fim, o Capítulo~\ref{sec:conclusao} apresenta as considerações finais do trabalho, discutindo as contribuições, limitações, atividades futuras e o cronograma de execução da pesquisa.
